\Needspace{11\baselineskip}
\SourceLeaf{273}
\begin{center}
\textit{Komposition der Geschlechter.}
\end{center}

\Article{246}
\textsc{Problem.} \textit{Es seien zwei beliebige primitive Formen $f,f'$ gegeben, durch deren Komposition $F$ entsteht. Aus den Geschlechtern, zu denen $f,f'$ gehören, ist das Geschlecht zu bestimmen, zu dem $F$ gehört.}

\textit{Lösung.} I. Wir betrachten zuerst den Fall, dass mindestens eine der Formen $f,f'$, etwa die erste, eigentlich primitiv ist, und bezeichnen die Determinanten der Formen $f,f',F$ mit $d,d',D$. Dann ist $D$ der größte gemeinsame Teiler von $dm'm'$, $d'$, wobei $m'$ entweder $1$ oder $2$ ist, je nachdem, ob die Form $f'$ eigentlich oder uneigentlich primitiv ist; $F$ gehört im ersten Fall zur eigentlich primitiven, im zweiten zur uneigentlich primitiven Ordnung. Das Geschlecht der Form $F$ wird durch ihre Einzelcharaktere bestimmt, und zwar sowohl bezüglich der einzelnen ungeraden Primteiler von $D$ als auch, in bestimmten Fällen, bezüglich der Zahlen $4$ oder $8$. Diese sind also einzeln zu bestimmen.

$1^{\circ}$. Ist $p$ irgendein ungerader Primteiler von $D$, so teilt diese Zahl notwendig auch $d,d'$; daher kommen unter den Charakteren der Formen $f,f'$ auch ihre Beziehungen zu $p$ vor. Kann nun $f$ die Zahl $a$ und $f'$ die Zahl $a'$ darstellen, so kann das Produkt $aa'$ durch $F$ dargestellt werden. Können also sowohl $f$ als auch $f'$ quadratische Reste modulo $p$ darstellen (die nicht durch $p$ teilbar sind), so kann auch $F$ quadratische Reste modulo $p$ darstellen; \textit{das heißt,}\ haben beide $f,f'$ den Charakter $Rp$, so hat die Form $F$ denselben Charakter. Ebenso hat $F$ den Charakter $Rp$, wenn beide $f,f'$ den Charakter $Np$ haben; dagegen hat $F$ den Charakter $Np$, wenn eine der Formen $f,f'$ den Charakter $Rp$, die andere $Np$ hat.

$2^{\circ}$. Geht in den Gesamtcharakter der Form $F$ eine Beziehung zur Zahl $4$ ein, so muss eine solche Beziehung auch in die Charaktere der Formen $f$, $f'$ eingehen. Denn das tritt nur ein, wenn $D$ entweder $\equiv0$ oder $\equiv3$ (modulo $4$) ist. Ist $D$ durch $4$ teilbar, so sind auch $dm'm'$ und $d'$ durch $4$ teilbar; daraus folgt sofort, dass $f'$ nicht uneigentlich primitiv sein kann, also $m'=1$. Somit sind sowohl $d$ als auch $d'$ durch $4$ teilbar, und in den Charakter beider geht eine Beziehung zu $4$ ein. Gilt $D\equiv3$ (modulo $4$), so teilt $D$ die Zahlen $d,d'$; die Quotienten sind Quadrate, also sind auch $d,d'$ notwendig entweder $\equiv0$ oder $\equiv3$ (modulo $4$), und unter den \SourceLeaf{274}Charakteren von $f,f'$ kommt eine Beziehung zu $4$ vor. Wie in ($1^{\circ}$) folgt daraus, dass die Form $F$ den Charakter $1,4$ hat, wenn entweder beide $f,f'$ den Charakter $1,4$ oder beide $3,4$ haben; dagegen hat die Form $F$ den Charakter $3,4$, wenn eine der Formen $f,f'$ den Charakter $1,4$, die andere $3,4$ hat.

$3^{\circ}$. Ist $D$ durch $8$ teilbar, so auch $d'$; daher ist $f'$ sicher eigentlich primitiv, $m'=1$, und auch $d$ ist durch $8$ teilbar. Somit kann unter den Charakteren der Form $F$ einer der Charaktere $1,8$; $3,8$; $5,8$; $7,8$ nur dann vorkommen, wenn sowohl im Charakter der Form $f$ als auch in dem der Form $f'$ eine solche Beziehung zu $8$ vorhanden ist. Wie zuvor bestätigt man leicht, dass die Form $F$ den Charakter $1,8$ hat, wenn $f$ und $f'$ bezüglich $8$ denselben Charakter haben; die Form $F$ hat den Charakter $3,8$, wenn eine der Formen $f,f'$ den Charakter $1,8$, die andere $3,8$ hat oder eine $5,8$, die andere $7,8$; $F$ hat $5,8$, wenn $f,f'$ die Charaktere $1,8$ und $5,8$ oder $3,8$ und $7,8$ haben; $F$ hat $.7,8$, wenn $f$ und $f'$ entweder $1,8$ und $7,8$ oder $3,8$ und $5,8$ haben.

$4^{\circ}$. Gilt $D\equiv2$ (modulo $8$), so ist $d'$ entweder $\equiv0$ oder $\equiv2$ (modulo $8$), also $m'=1$ und somit auch $d$ entweder $\equiv0$ oder $\equiv2$ (modulo $8$). Allerdings können nicht beide $d,d'$ durch $8$ teilbar sein, da $D$ ihr \textit{größter} gemeinsamer Teiler ist. Daher kann einer der Charaktere $1\text{ \textit{und} }7,8$; $3\text{ \textit{und} }5,8$ der Form $F$ nur dann zukommen, wenn entweder beide Formen $f,f'$ einen davon haben oder eine einen davon und die andere einen der Charaktere $1,8$; $3,8$; $5,8$; $7,8$. Daraus folgt leicht, dass der Charakter der Form $F$ durch die folgende Tabelle bestimmt wird, wenn der Charakter am Rand zu einer der Formen $f,f'$ gehört, der oben stehende Charakter zur anderen:
\[
\begin{array}{||c||c|c||}
\hline\hline
 & 1\text{ \textit{und} }7,8 & 3\text{ \textit{und} }5,8 \\
 & \text{oder }1,8 & \text{oder }3,8 \\
 & \text{oder }7,8 & \text{oder }5,8 \\
\hline\hline
1\text{ \textit{und} }7,8 & 1\text{ \textit{und} }7,8 & 3\text{ \textit{und} }5,8 \\
\hline
3\text{ \textit{und} }5,8 & 3\text{ \textit{und} }5,8 & 1\text{ \textit{und} }7.8 \\
\hline\hline
\end{array}
\]

$5^{\circ}$. Ebenso beweist man, dass der Form $F$ keiner der Charaktere $1\text{ \textit{und} }3,8$; $5\text{ \textit{und} }7,8$ zukommen kann, wenn nicht auch mindestens eine der Formen $f,f'$ einen davon hat und die andere entweder einen davon oder einen der Charaktere $1,8$; $3,8$; $5,8$; $7,8$.
\SourceLeaf{275}
Dabei wird der Charakter der Form $F$ durch diese Tabelle bestimmt, an deren Rand und oben die Charaktere der Formen $f,f'$ stehen:
\[
\begin{array}{||c||c|c||}
\hline\hline
 & 1\text{ \textit{und} }3,8 & 5\text{ \textit{und} }7,8 \\
 & \text{oder }1,8 & \text{oder }5,8 \\
 & \text{oder }3,8 & \text{oder }7,8 \\
\hline\hline
1\text{ \textit{und} }3,8 & 1\text{ \textit{und} }3,8 & 5\text{ \textit{und} }7,8 \\
\hline
5\text{ \textit{und} }7,8 & 5\text{ \textit{und} }7,8 & 1\text{ \textit{und} }3.8 \\
\hline\hline
\end{array}
\]

II. Sind beide Formen $f,f'$ uneigentlich primitiv, so ist $D$ der größte gemeinsame Teiler von $4d.\ 4d'$, oder $\frac{1}{4}D$ der größte gemeinsame Teiler von $d.\ d'$. Daraus folgt leicht, dass sowohl $d$ als auch $d'$ und $\frac{1}{4}D$ jeweils $\equiv1$ (modulo $4$) sind. Setzt man $F=(A,B,C)$, so ist der größte gemeinsame Teiler von $A,B,C$ $=2$. und der größte gemeinsame Teiler von $A,2B,C$ gleich $4$. Daher ist $F$ eine von der uneigentlich primitiven Form $(\frac{1}{2}A,\frac{1}{2}B,\frac{1}{2}C)$ abgeleitete Form; deren Determinante ist $\frac{1}{4}D$, und ihr Geschlecht bestimmt das Geschlecht der Form $F$. Da jene Form uneigentlich primitiv ist, enthält ihr Charakter keine Beziehungen zu $4$ oder $8$, sondern nur zu den einzelnen ungeraden Primteilern von $\frac{1}{4}D$. Alle diese Teiler teilen nun offenbar auch $d,d'$, und die Hälfte jedes Produkts zweier Faktoren. deren einer durch $f$, der andere durch $f'$ darstellbar ist, kann durch die Form $(\frac{1}{2}A.\ \frac{1}{2}B,\frac{1}{2}C)$ dargestellt werden. Daher erkennt man leicht, dass der Charakter dieser Form bezüglich jeder ungeraden Primzahl $p$, die $\frac{1}{4}D$ teilt, gleich $Rp$ ist, \textit{sowohl} wenn $2Rp$ gilt und die Formen $f,f'$ bezüglich $p$ denselben Charakter haben. \textit{als auch} wenn $2Np$ gilt und die Charaktere der Formen $f,f'$ bezüglich $p$ entgegengesetzt sind; dagegen ist der Charakter jener Form $Np$, \textit{sowohl} wenn $f,f'$ bezüglich $p$ gleiche Charaktere haben und $2Np$ gilt, \textit{als auch} wenn $f,f'$ entgegengesetzte Charaktere haben und $2Rp$ gilt.

\Needspace{8\baselineskip}
\Article{247}
Aus der Lösung des vorigen Problems ist klar: Ist $g$ eine primitive Form derselben Ordnung und desselben Geschlechts wie $f$ und ebenso $g'$ eine primitive Form derselben Ordnung und desselben Geschlechts wie $f'$, so gehört die aus $g$ und $g'$ zusammengesetzte Form zu demselben Geschlecht wie die aus $f$ und $f'$ zusammengesetzte. Daraus ergibt sich die Bedeutung eines \textit{Geschlechts},
\SourceLeaf{276}
das aus zwei anderen Geschlechtern (oder auch mehreren) \textit{zusammengesetzt} ist. Ferner folgt daraus: Haben $f,f'$ dieselbe Determinante, ist $f$ eine Form aus dem Hauptgeschlecht und $F$ aus $f$ und $f'$ zusammengesetzt, so gehört $F$ zu demselben Geschlecht wie $f'$; daher kann das Hauptgeschlecht bei der Komposition mit anderen Geschlechtern derselben Determinante stets weggelassen werden. Gehört dagegen bei sonst gleichen Bedingungen $f$ nicht zum Hauptgeschlecht, während $f'$ eine primitive Form ist, so gehört $F$ sicher zu einem anderen Geschlecht als $f'$. Sind schließlich $f,f'$ eigentlich primitive Formen desselben Geschlechts, so gehört $F$ zum Hauptgeschlecht; sind dagegen $f,f'$ beide eigentlich primitiv und von derselben Determinante, gehören aber zu verschiedenen Geschlechtern, so kann $F$ nicht zum Hauptgeschlecht gehören. Wird also irgendeine eigentlich primitive Form \textit{mit sich selbst} zusammengesetzt, so gehört die entstehende Form, die ebenfalls eigentlich primitiv und von derselben Determinante ist, notwendig zum Hauptgeschlecht.

\Needspace{8\baselineskip}
\Article{248}
\textsc{Problem.} \textit{Es seien zwei beliebige Formen $f,f'$ gegeben, aus denen $F$ zusammengesetzt ist. Aus den Geschlechtern der Formen $f,f'$ ist das Geschlecht der Form $F$ zu bestimmen.}

\textit{Lösung.} Es sei $f=(a,b,c)$, $f'=(a',b',c')$, $F=(A,B,C)$, ferner $\mathfrak{m}$ der größte gemeinsame Teiler von $a,b,c$ und $\mathfrak{m}'$ der größte gemeinsame Teiler von $a',b',c'$, sodass $f,f'$ von den primitiven Formen $(\frac{a}{\mathfrak{m}},\allowbreak \frac{b}{\mathfrak{m}},\allowbreak \frac{c}{\mathfrak{m}})$, $(\frac{a'}{\mathfrak{m}'},\allowbreak \frac{b'}{\mathfrak{m}'},\allowbreak \frac{c'}{\mathfrak{m}'})$ abgeleitet sind, die wir jeweils mit $\mathfrak{f}$, $\mathfrak{f}'$ bezeichnen. Ist nun mindestens eine der Formen $\mathfrak{f}$, $\mathfrak{f}'$ eigentlich primitiv, so ist der größte gemeinsame Teiler von $A,B,C$ gleich $\mathfrak{m}\mathfrak{m}'$ und daher $F$ von der primitiven Form $(\frac{A}{\mathfrak{m}\mathfrak{m}'},\allowbreak \frac{B}{\mathfrak{m}\mathfrak{m}'},\allowbreak \frac{C}{\mathfrak{m}\mathfrak{m}'})\ldots\mathfrak{F}$ abgeleitet; somit hängt das Geschlecht der Form $F$ vom Geschlecht der Form $\mathfrak{F}$ ab. Man erkennt aber leicht, dass $\mathfrak{F}$ durch dieselbe Substitution in $\mathfrak{f}\mathfrak{f}'$ übergeht, durch die $F$ in $ff'$ übergeht. Also ist $\mathfrak{F}$ aus $\mathfrak{f}$, $\mathfrak{f}'$ zusammengesetzt, und ihr Geschlecht kann durch das Problem des Artikels~246 bestimmt werden. --- Sind dagegen beide $\mathfrak{f}$, $\mathfrak{f}'$ uneigentlich primitiv, so ist der größte gemeinsame Teiler von $A,B,C$ gleich $2\mathfrak{m}\mathfrak{m}'$; die Form $\mathfrak{F}$ ist weiterhin aus $\mathfrak{f}$, $\mathfrak{f}'$ zusammengesetzt und offenbar von der eigentlich primitiven Form $(\frac{A}{2\mathfrak{m}\mathfrak{m}'},\allowbreak \frac{B}{2\mathfrak{m}\mathfrak{m}'},\allowbreak \frac{C}{2\mathfrak{m}\mathfrak{m}'})$ abgeleitet. Das Geschlecht dieser Form kann daher nach Artikel~246 bestimmt werden; da $F$ von derselben Form abgeleitet ist, ergibt sich daraus sofort auch ihr Geschlecht.

Aus dieser Lösung ist klar, dass der im vorigen Artikel für primitive Formen erläuterte Satz, nämlich \textit{gehören $f',g'$ jeweils zu denselben Geschlechtern wie $f,g$, so gehört die aus $f',g'$ zusammengesetzte Form zu demselben Geschlecht wie die aus $f,g$ zusammengesetzte}, allgemein für beliebige Formen gilt.
